% Les emplacements d'image du rapport, et les seules commandes par lesquelles une
% image entre dans le document.
%
% DEUX FAMILLES D'IMAGE SONT ADMISES, ET DEUX SEULEMENT : le logotype de
% l'établissement de formation, et les captures du tableau de bord PRODUIT PAR CE
% PROJET. Toute image du système d'information observé reste interdite, sans
% exception — c'est la décision `0010`, que rien ici n'assouplit.
% `tests/test_aucune_image.py` tient la propriété sur les fichiers suivis : seuls
% `report/figures/logos/` et `report/figures/tableau-de-bord/` la portent.
%
% LA PLACE EST RÉSERVÉE MAINTENANT ET REMPLIE PLUS TARD. Chaque commande teste
% l'existence du fichier attendu. S'il existe, elle l'inclut. S'il n'existe pas,
% elle compose un cadre de la MÊME hauteur ou de la même largeur déclarée, portant
% la légende et la mention « à insérer ». Le document compose donc à tous les
% stades, et la pagination ne bouge presque pas à l'insertion — c'est tout
% l'intérêt du mécanisme, et une commande qui composerait un cadre d'une autre
% taille que l'image finale l'annulerait.
%
% L'ÉCART RÉSIDUEL EST MESURÉ, PAS SUPPOSÉ, ET LA MESURE A CORRIGÉ LE CODE. Le
% cadre est posé avec `\fboxsep` nul et aligné PAR LE BAS — `\parbox[b]`, comme
% `\includegraphics`, dont la ligne de base est le bas de l'image. Un premier
% essai alignait le cadre sur son centre (`\parbox[c]`) : un document témoin
% composé avec puis sans l'image montrait alors la ligne suivante déplacée de
% 3,304 pt. Avec l'alignement par le bas, la même mesure donne 0,395 pt — la seule
% épaisseur du trait inférieur du cadre, soit 0,047 % de la hauteur utile d'une
% page A4. Un emplacement rempli déplace donc le texte qui le suit d'un vingtième
% de ligne, et la pagination ne bouge pas.
%
% La mesure : composer un document portant `\captureTdb` et un mot repérable
% après lui, relever `yMin` de ce mot par `pdftotext -bbox`, avec le fichier image
% en place puis retiré. Les deux relevés sont au rapport de ce travail.

% Le trait du cadre d'attente, et l'absence de marge intérieure. Les deux sont
% posés localement par chaque commande, jamais globalement : `\fbox` sert ailleurs
% dans le document — la mention de version de travail en porte un — et lui retirer
% sa marge partout changerait cette mention-là aussi.
\newcommand{\cadreDAttenteReglages}{\setlength{\fboxsep}{0pt}\setlength{\fboxrule}{0.4pt}}

% Deux longueurs de travail. Les dimensions ne se calculent pas par juxtaposition
% dans un argument de macro : `4#1` avec `#1` valant `2.2cm` donnerait `42.2cm`, et
% `0.625#2` produirait une dimension inexistante. Le calcul passe donc par
% `\setlength`, qui admet un facteur devant une longueur.
\newlength{\largeurEmplacement}
\newlength{\hauteurEmplacement}

% Un cadre d'attente de dimensions imposées. La largeur et la hauteur sont celles
% que l'image occupera ; le texte est centré dans les deux sens et n'influe pas
% sur la boîte, `\parbox[c][hauteur][c]` fixant la hauteur indépendamment de son
% contenu.
\newcommand{\cadreDAttente}[3]{%
  {\cadreDAttenteReglages%
   \fbox{\parbox[b][#2][c]{#1}{\centering\footnotesize\itshape #3\par}}}%
}

% ------------------------------------------------------------------------------
% LES DEUX LOGOTYPES INSTITUTIONNELS.
%
%   \logosInstitutionnels
%
% La page de garde en porte DEUX, côte à côte : l'établissement de formation et
% l'organisme dont il relève. La commande précédente, `\logoEcole{hauteur}`, n'en
% prenait qu'un et cherchait un fichier `ecole.pdf` ou `ecole.png` qui n'a jamais
% existé ; elle est remplacée, et non doublée, pour qu'il n'y ait qu'une seule
% voie par laquelle un logotype entre dans le document.
%
% LES DEUX HAUTEURS NE SONT PAS ÉGALES, ET C'EST MESURÉ. Imposer la même hauteur
% aux deux ne les équilibre pas : un logotype large et bas paraît plus grand qu'un
% logotype compact et haut à hauteur égale. L'égalisation porte donc sur la
% SURFACE OPTIQUE — la surface de la boîte utile, marge transparente exclue.
%
%   fichier          dimensions   rapport   boîte utile   marge   encre        hauteur
%   INSEA-logo.png   1200x1304    0,9202    983x1110      30,3 %  582 421 px   20,0 mm
%   HCP-logo.png     3790x1088    3,4835    3781x1088      0,2 %  547 086 px   17,2 mm
%
% Le calcul : la surface portée à la page vaut la surface en pixels multipliée par
% le carré de l'échelle, et l'échelle vaut la hauteur déclarée divisée par la
% hauteur du fichier. Égaliser donne un rapport d'échelles égal à la racine du
% rapport des surfaces, et le rapport des hauteurs s'en déduit en multipliant par
% le rapport des hauteurs en pixels.
%
% DEUX SURFACES SONT CANDIDATES, ET ELLES NE DONNENT PAS LE MÊME RÉSULTAT. La
% boîte utile donne un rapport de hauteurs de 2,3267, donc 8,6 mm pour le second
% logotype ; l'encre — le nombre de pixels réellement opaques — donne 1,1616, donc
% 17,2 mm. L'écart tient à ceci : la boîte utile de l'INSEA est couverte à 53,4 %,
% celle du HCP à 13,3 % seulement, parce que ce second logotype est un mot écrit
% en fin, largement étalé, quand le premier est une marque pleine. La boîte utile
% créditait donc le HCP d'une surface qu'il n'occupe pas.
%
% LA MESURE À L'ENCRE EST RETENUE, ET LE RENDU L'A TRANCHÉE. La page composée avec
% 8,6 mm a été regardée en image : le second logotype y paraissait nettement plus
% faible que le premier, et sa ligne « Haut-Commissariat au Plan » n'était plus
% lisible. À 17,2 mm les deux marques pèsent le même poids à l'œil et les deux
% textes se lisent. C'est la surface d'encre qui mesure la présence optique, pas
% le rectangle qui l'entoure.
%
% LES DEUX SONT ALIGNÉS SUR LEUR LIGNE MÉDIANE, jamais sur leur base : deux
% logotypes de hauteurs différentes alignés par le bas paraissent décalés.
% `\raisebox{-0.5\height}` pose chaque boîte à cheval sur la ligne de base, si
% bien que les deux milieux se retrouvent à la même hauteur. L'encombrement
% vertical de la ligne reste celui du plus grand des deux, et il ne dépend pas de
% la présence des fichiers : les cadres d'attente ont les mêmes dimensions.
%
% CENTRER LES BOÎTES NE CENTRE PAS L'ENCRE, et la mesure l'a montré. Le fichier de
% l'INSEA porte 51 pixels de marge transparente en haut et 143 en bas : son encre
% est donc plus haute que le centre de sa boîte de 46 pixels, soit 0,706 mm à la
% hauteur déclarée. Le fichier du HCP, lui, n'a aucune marge verticale. Les deux
% boîtes alignées, les deux encres l'étaient à 0,76 mm près — relevé sur l'image du
% PDF, et cohérent avec le calcul. Chaque logotype porte donc une CORRECTION,
% mesurée sur son fichier, qui ramène sa ligne médiane d'encre sur la ligne de
% base. La correction s'applique au `\raisebox` extérieur, donc aussi bien à
% l'image qu'à son cadre d'attente.
\newcommand{\hauteurLogoEtablissement}{20mm}
\newcommand{\hauteurLogoTutelle}{17.2mm}
\newcommand{\ecartEntreLogos}{16mm}

% Un logotype, ou son cadre d'attente aux mêmes dimensions.
%
%   \unLogo{base du nom}{hauteur}{rapport largeur/hauteur}{correction}{légende}
%
% Le rapport est celui du FICHIER, mesuré sur lui, et non une valeur commode : il
% ne sert qu'au cadre d'attente, dont la largeur doit être celle qu'occupera
% l'image pour que retirer un fichier ne déplace rien, ni verticalement ni
% horizontalement. La correction est la longueur dont il faut descendre la boîte
% pour que sa médiane d'encre tombe sur la ligne de base ; elle vaut zéro pour un
% fichier sans marge transparente.
\newcommand{\unLogo}[5]{%
  \raisebox{\dimexpr-0.5\height-#4\relax}{%
    \IfFileExists{figures/logos/#1.pdf}{%
      \includegraphics[height=#2]{figures/logos/#1.pdf}%
    }{%
      \IfFileExists{figures/logos/#1.png}{%
        \includegraphics[height=#2]{figures/logos/#1.png}%
      }{%
        \setlength{\hauteurEmplacement}{#2}%
        \setlength{\largeurEmplacement}{#3\hauteurEmplacement}%
        \cadreDAttente{\largeurEmplacement}{\hauteurEmplacement}{#5\\à insérer}%
      }%
    }%
  }%
}

\newcommand{\logosInstitutionnels}{%
  \unLogo{INSEA-logo}{\hauteurLogoEtablissement}{0.9202}{0.706mm}{logotype de l'établissement}%
  \hspace{\ecartEntreLogos}%
  \unLogo{HCP-logo}{\hauteurLogoTutelle}{3.4835}{0mm}{logotype de la tutelle}%
}

% ------------------------------------------------------------------------------
% Une capture du tableau de bord produit par ce projet.
%
%   \captureTdb{fichier}{largeur}{légende}{label}
%
% `fichier` est le nom seul, sans répertoire ni extension : la commande le cherche
% sous `figures/tableau-de-bord/`, en `.png` puis en `.pdf`. Écrire le répertoire à
% chaque appel l'exposerait à diverger du préfixe que le contrôle tolère.
%
% LA HAUTEUR RÉSERVÉE VIENT DE LA LARGEUR, PAR UN RAPPORT DÉCLARÉ PAR CAPTURE.
% Une capture n'a pas de hauteur déclarée à l'appel, et le cadre d'attente doit
% pourtant en avoir une. Un premier état retenait un rapport UNIQUE de 16:10 pour
% toutes les captures, en écrivant que c'était le point aveugle du mécanisme :
% « une capture d'un autre rapport se compose correctement, mais déplace ce qui la
% suit ». Les trois captures déposées ont démenti la supposition — 2,456, 1,973 et
% 1,964 au lieu de 1,600 —, et le rapport unique réservait donc de 1,7 à 3,2 cm de
% trop pour chacune. Le rapport est désormais DÉCLARÉ POUR CHAQUE FICHIER, mesuré
% sur le fichier lui-même ; la valeur unique ne subsiste que comme défaut, pour une
% capture non encore déposée et donc non encore mesurable.
%
% La déclaration est ici, dans le fichier qui tient les images, et non aux appels :
% c'est une propriété du fichier image, pas du texte qui l'appelle.
\newcommand{\rapportHauteurCapture}{0.625}

\newcommand{\declarerCapture}[2]{%
  \expandafter\gdef\csname rapportCapture@#1\endcsname{#2}%
}
\newcommand{\rapportDeLaCapture}[1]{%
  \ifcsname rapportCapture@#1\endcsname
    \csname rapportCapture@#1\endcsname
  \else
    \rapportHauteurCapture
  \fi
}

% Hauteur divisée par largeur, en pixels, de chaque fichier déposé.
\declarerCapture{page-activite}{0.40724}%  1374/3374
\declarerCapture{page-donnees}{0.50680}%   1714/3382
\declarerCapture{page-sejours}{0.50918}%   1720/3378

% La numérotation. Le rapport ne porte aujourd'hui aucun environnement flottant et
% aucune `\caption` ; un `\label` posé sans compteur renverrait à la section
% courante, ce qui ferait mentir tout `\ref`. Un compteur propre est donc déclaré,
% et `\refstepcounter` en fait la cible du `\label` qui suit.
\newcounter{capturetdb}
\renewcommand{\thecapturetdb}{\arabic{capturetdb}}

\newcommand{\legendeCapture}[2]{%
  \begin{center}%
    \footnotesize Figure~\thecapturetdb~— #1\label{#2}%
  \end{center}%
}

% L'IMAGE EST POSÉE DANS LA BOÎTE RÉSERVÉE, ET NON À CÔTÉ. La boîte a la hauteur
% déclarée pour ce fichier ; l'image y entre bornée en largeur ET en hauteur, à
% rapport conservé. Un fichier dont le rapport dériverait de sa déclaration ne
% déborderait donc pas de la réservation : il y laisserait un peu de blanc. C'est
% ce qui rend la pagination indépendante de la présence du fichier, au lieu de la
% faire dépendre d'une supposition écrite ailleurs.
\newcommand{\captureTdb}[4]{%
  \refstepcounter{capturetdb}%
  \setlength{\largeurEmplacement}{#2}%
  \edef\facteurDeLaCapture{\rapportDeLaCapture{#1}}%
  \setlength{\hauteurEmplacement}{\facteurDeLaCapture\largeurEmplacement}%
  \begin{center}%
    \IfFileExists{figures/tableau-de-bord/#1.png}{%
      \parbox[b][\hauteurEmplacement][c]{\largeurEmplacement}{\centering%
        \includegraphics[width=\largeurEmplacement,height=\hauteurEmplacement,keepaspectratio]{figures/tableau-de-bord/#1.png}}%
    }{%
      \IfFileExists{figures/tableau-de-bord/#1.pdf}{%
        \parbox[b][\hauteurEmplacement][c]{\largeurEmplacement}{\centering%
          \includegraphics[width=\largeurEmplacement,height=\hauteurEmplacement,keepaspectratio]{figures/tableau-de-bord/#1.pdf}}%
      }{%
        \cadreDAttente{\largeurEmplacement}{\hauteurEmplacement}{#3\\capture à insérer}%
      }%
    }%
  \end{center}%
  \legendeCapture{#3}{#4}%
}
